% =====================================================================
%  Currículo — Augusto de Castro Menchaca
%  Estrutura seguindo o CV de Amanda Vieira (Product Designer).
%
%  Compilar com pdfLaTeX. Nada além de pacotes padrão do TeX Live/MiKTeX.
%  Sem LaTeX instalado? Suba este arquivo no overleaf.com e compile lá.
%
%  ATENÇÃO: tudo marcado com \pend{...} sai em VERMELHO no PDF.
%  São dados que eu não tinha e que NÃO inventei. Preencha e apague a macro.
% =====================================================================
\documentclass[11pt,a4paper]{article}

\usepackage[utf8]{inputenc}
\usepackage[T1]{fontenc}
\usepackage[brazil]{babel}
\usepackage[margin=1.8cm]{geometry}
\usepackage{helvet}
\renewcommand{\familydefault}{\sfdefault}
\usepackage{titlesec}
\usepackage{enumitem}
\usepackage{xcolor}
\usepackage[hidelinks]{hyperref}
\usepackage{microtype}

% ---------- cores ----------
\definecolor{tinta}{HTML}{1A1A1A}
\definecolor{cinza}{HTML}{5A5A5A}
\definecolor{falta}{HTML}{B3261E}
\AtBeginDocument{\color{tinta}}   % \color solto no preâmbulo não compila

% ---------- marcador de dado pendente ----------
\newcommand{\pend}[1]{\textcolor{falta}{\textbf{[#1]}}}

% ---------- seções com régua, como no CV de referência ----------
\titleformat{\section}
  {\large\bfseries}{}{0pt}{}[\vspace{2pt}\titlerule\vspace{-2pt}]
\titlespacing*{\section}{0pt}{11pt}{6pt}

% ---------- entrada de experiência ----------
% \entrada{Organização, Cargo}{período}
\newcommand{\entrada}[2]{%
  \vspace{4pt}%
  \noindent\textbf{#1}\\[1pt]%
  {\small\itshape\color{cinza} #2}\par\vspace{2pt}%
}

\setlist[itemize]{leftmargin=14pt,itemsep=1.5pt,topsep=2pt,parsep=0pt}
\setlength{\parindent}{0pt}
\pagestyle{empty}

% =====================================================================
\begin{document}

% ------------------------------- CABEÇALHO
{\LARGE\bfseries Augusto de Castro Menchaca}\\[3pt]
{\large\bfseries Desenvolvimento de Software \textbar{} Computação Científica \textbar{} Gestão de Projetos}\\[4pt]
{\small\color{cinza}%
Pelotas, RS, Brasil \textbar{}
(53) 99901-1310 \textbar{}
\href{mailto:adcmenchaca@inf.ufpel.edu.br}{adcmenchaca@inf.ufpel.edu.br}\\
\href{https://linkedin.com/in/augusto-menchaca-078469330}{linkedin.com/in/augusto-menchaca-078469330} \textbar{}
\href{https://github.com/AugustoMenchaca}{github.com/AugustoMenchaca} \textbar{}
\href{https://augustomenchaca.github.io/}{augustomenchaca.github.io}%
}

% ------------------------------- PERFIL
\section*{Perfil Profissional}
Graduando em Ciência da Computação na UFPel, atuando entre desenvolvimento de
software, computação científica e condução de projetos de tecnologia. No Núcleo
Integrado de Previsão (NIP/UFPel), como bolsista, respondo pela camada de
software de ferramentas científicas e trabalho na aplicação de machine learning
a recursos hídricos. Na Hut 8, empresa júnior de computação da UFPel, sou
Diretor de Projetos: conduzo projetos de clientes do discovery à entrega,
coordenando requisitos, escopo, squads, QA e infraestrutura. Meu interesse está
no trecho entre um problema pouco definido e uma solução que continua
funcionando depois que eu saio.

% ------------------------------- EXPERIÊNCIA
\section*{Experiência}

\entrada{Núcleo Integrado de Previsão (NIP) --- UFPel, Pesquisador de Graduação (Bolsista)}
        {set. 2025 --- Presente \textbar{} Pelotas, RS}
\begin{itemize}
  \item \textbf{Camada de software do IDF-BR:} responsável pela escolha de
        ferramentas, decisões técnicas, arquitetura, estrutura da aplicação,
        fluxo, UI/UX, implementação e deploy da plataforma
        \href{https://www.idf-br.com.br/home/index.html}{idf-br.com.br}. O time
        científico definiu dados, equações e outputs; a camada de software foi
        minha responsabilidade.
  \item \textbf{Estruturação de dados científicos:} transformação do acervo de
        planilhas do grupo de pesquisa em JSON estruturado, tornando os dados
        consumíveis por uma aplicação web.
  \item \textbf{IA aplicada a recursos hídricos:} escola de verão \textbf{HydroUAI},
        na UFMG, em janeiro de 2026, a partir da qual o foco do trabalho passou
        a se aproximar de machine learning, experimentação e escrita
        acadêmica.
  \item \textbf{Pesquisa aplicada:} experimentação em machine learning voltada a
        recursos hídricos, com o objetivo de transformar pesquisa complexa em
        ferramentas acessíveis fora da universidade.
\end{itemize}

\entrada{Hut 8 --- Empresa Júnior de Computação da UFPel, Diretor de Projetos}
        {nov. 2024 --- Presente \textbar{} Diretoria desde jun. 2025 \textbar{} Pelotas, RS}
\begin{itemize}
  \item \textbf{Condução de projetos de cliente:} atuação do discovery à
        entrega, incluindo levantamento de requisitos, definição de escopo,
        precificação e elaboração de proposta comercial.
  \item \textbf{Coordenação de squads:} formação de times, acompanhamento de
        execução, apoio ao gerente de projeto e ponte entre design e
        desenvolvimento.
  \item \textbf{QA e infraestrutura:} condução de QA com o cliente, configuração
        de banco de dados, bucket/armazenamento e deploy de API e frontend.
  \item \textbf{Responsabilidade organizacional:} assumi a diretoria em um
        período de baixa atividade da empresa júnior, quando a continuidade da
        operação estava ameaçada, aprendendo gestão na prática.
  \item \textbf{DVO --- Prefeitura de Pelotas}
        (\href{https://dvopelotas.com.br/}{dvopelotas.com.br}): coordenação do
        projeto dos requisitos à produção. O sistema foi construído pelo squad;
        meu papel foi coordenação, escopo e apoio à entrega, além da
        configuração de banco, armazenamento e dos deploys.
  \item \textbf{Ciere da Rosa --- Advocacia de Família e Sucessões}
        (\href{https://advocaciacieredarosa.com.br/}{advocaciacieredarosa.com.br}):
        atuação comercial, discovery, requisitos e escopo, sendo a ponte entre
        cliente, design e desenvolvimento. Havia liderança de design própria no
        projeto. Responsável também por QA, SEO e pela atenção a GEO --- como o
        escritório é descrito e recuperado por modelos de linguagem, não só por
        buscadores.
\end{itemize}

% ------------------------------- EDUCAÇÃO
\section*{Educação}

\entrada{Bacharelado em Ciência da Computação, Universidade Federal de Pelotas (UFPel)}
        {2024/1 --- previsão de conclusão entre 2027/2 e 2028/1}

% ------------------------------- PROJETOS
\section*{Projetos}

\entrada{Quantum Machine Learning --- Pesquisa Experimental}
        {jul. --- ago. 2026 \textbar{} Disciplina da graduação \textbar{} Projeto acadêmico, sem publicação}
\begin{itemize}
  \item \textbf{Pergunta central:} o QML tem desempenho melhor simplesmente por
        usar representações quânticas, ou o desempenho depende do alinhamento
        entre dados, feature maps e kernels?
  \item \textbf{Desenho experimental:} comparação controlada entre SVM clássico,
        QSVC e métodos de kernel quântico, sobre dados reais (Breast Cancer
        Wisconsin) e sintéticos (Qiskit \texttt{ad\_hoc\_data}).
  \item \textbf{Resultado:} nos dados reais, o SVM clássico e a configuração
        quântica alcançaram resultados comparáveis na representação controlada,
        sem vantagem quântica clara. No benchmark sintético o modelo quântico
        teve desempenho excelente, mas o conjunto de dados havia sido desenhado
        em torno de uma representação quântica compatível. A conclusão é que o
        resultado depende da combinação entre dados, representação, feature map
        e kernel --- não do uso de computação quântica em si.
\end{itemize}

\entrada{Inteligência Artificial aplicada à Engenharia Hídrica --- Predição de Fluxo}
        {Projeto acadêmico em andamento --- NIP/UFPel}
\begin{itemize}
  \item \textbf{Objetivo:} aplicar métodos de machine learning à predição de
        fluxo em sistemas hídricos, dentro da linha de pesquisa do Núcleo
        Integrado de Previsão.
  \item \textbf{Relevância:} previsão de fluxo é insumo direto para alerta e
        planejamento em bacias hidrográficas. É pesquisa acadêmica com aplicação
        social clara, na direção de tornar previsão hidrológica acessível fora
        da universidade.
\end{itemize}

% ------------------------------- CERTIFICAÇÕES
\section*{Certificações \& Cursos}
\begin{itemize}
  \item \textbf{HydroUAI --- Escola de Verão em Inteligência Artificial aplicada a
        Recursos Hídricos:} Universidade Federal de Minas Gerais (UFMG),
        19 a 23 de janeiro de 2026. Previsão de vazão afluente com redes
        neurais (MLP, LSTM, GRU e mecanismos de atenção), preenchimento de
        falhas em séries hidrológicas, seleção de atributos, construção de
        defasagens e busca de hiperparâmetros.
\end{itemize}

% ------------------------------- COMPETÊNCIAS
\section*{Competências}

\textbf{Linguagens:} Python, C, C++, Go, JavaScript, HTML, CSS --- experiência
de uso em contexto acadêmico e em projetos próprios

\vspace{3pt}
\textbf{Frameworks \& Bibliotecas:} PyTorch, scikit-learn, pandas, NumPy, SciPy,
Matplotlib e Seaborn; GeoPandas, Rasterio, Shapely, PyProj, Folium e Google Earth
Engine para dados geoespaciais

\vspace{3pt}
\textbf{Dados \& Infraestrutura:} estruturação de dados científicos
(planilha $\rightarrow$ JSON), configuração de banco de dados,
bucket/armazenamento, deploy de API e de frontend

\vspace{3pt}
\textbf{Gestão de Projetos:} levantamento de requisitos, definição de escopo,
precificação, proposta comercial, formação e coordenação de squad, QA com
cliente, ponte design--desenvolvimento, acompanhamento de entrega

\vspace{3pt}
\textbf{Pesquisa:} computação científica, machine learning aplicado, recursos
hídricos, experimentação controlada, escrita acadêmica

\vspace{3pt}
\textbf{Idiomas:} Português --- nativo; Inglês --- intermediário

\end{document}
